\documentclass[11pt]{article}

\usepackage[margin=1in]{geometry}
\usepackage{kotex}
\usepackage[T1]{fontenc}
\usepackage[utf8]{inputenc}
\usepackage{lmodern}
\usepackage{microtype}
\usepackage{amsmath,amssymb,mathtools}
\usepackage{booktabs}
\usepackage{array}
\usepackage{tabularx}
\usepackage{hyperref}

\hypersetup{
  colorlinks=true,
  linkcolor=blue!55!black,
  urlcolor=blue!55!black,
  citecolor=blue!55!black
}

\title{MCF v6 Lower Bound Model with Canal Detours and Relaxed TS Timing}
\author{Yongsu Rah}
\date{2026년 5월 8일}

\begin{document}
\maketitle

\begin{abstract}
이 문서는 \texttt{algorithms/yongs/mcf\_v6\_lb}에서 \texttt{mcf\_v5} 대비
달라지는 부분만 기술한다. 공통적인 multicommodity flow 네트워크, position
activation, service coverage, TS activation/balance, dry dock 보정, 목적함수의
기본 구조는 \texttt{mcf\_v5}와 동일하다. \texttt{mcf\_v6\_lb}의 목적은 TS 발생
시각이 고정되지 않은 정확한 원 문제의 최적값에 대한 lower bound를 얻는 것이다.
이를 위해 TS 관련 항해 가능성을 낙관적으로 평가하는 완화 네트워크를 구성하고,
TS 및 TS 관련 항해 비용은 가능한 정확 문제 비용의 lower envelope로 계산한다.
항로 밖 운하 경유 후보는 \texttt{CanalSailArc}로 유지하고, 운하 route bunker 비용도
TS 관련 arc인 경우 동일한 lower-envelope 정책으로 계산한다. Canal fee는 입력상
연월 기준이 없는 고정 비용이므로 lower envelope를 적용하지 않는다.
운하 경유 가능 여부는 \texttt{mcf\_v5}와 동일하게
\texttt{cost\_canal\_direction.csv}의 nonempty direction row로 정의한다.
\end{abstract}

\section{목적}

정확한 원 문제에서는 TS 발생 시각이 고정되어 있지 않다. 반면 기존 MCF 계열 모형은
TS 발생 시각을 정규 port stay의 직전 또는 직후로 고정한다(TS restriction). 이 제한 때문에
기존 MCF 최적값은 정확한 원 문제의 최적값과 다를 수 있다.

\texttt{mcf\_v6\_lb}는 정확한 원 문제를 직접 모델링하지 않는다. 대신 정확한 원
문제에서 가능한 TS timing을 포함하도록 항해 시간 제약을 완화한 문제를 만든다.
다만 현재 구현의 기본 설정은 실제 선박만 사용하는 모형이다. 따라서 lower bound
statement는 가상 선박을 포함한 전체 feasible set inclusion으로 쓰지 않고, 다음
가정하의 value inequality로만 해석한다.

\begin{itemize}
  \item TS time restriction 문제와 free-TS exact 문제 모두에서 가상 선박을 사용하지
  않는 최적해가 존재한다.
  \item 완화된 TS 시간으로 계산한 항해 가능 시간은 exact timing에서 필요한 항해
  가능 시간보다 짧지 않으며, 항해 가능 시간이 길어질수록 필요한 속도와 bunker
  consumption은 증가하지 않는다.
  \item TS activation cost와 TS 관련 항해의 bunker price month attribution은 exact
  가능한 비용의 lower envelope로 계산한다.
  \item $z^{LB}$와 비교되는 값은 같은 MIP objective 기준이다. OCAM
  evaluation total cost 기준의 certified bound라고 주장하려면 별도의 증명이 필요하다.
\end{itemize}

위 가정이 성립하면 minimization 기준에서
\[
  z^{LB}_{\mathrm{MIP}} \le z^{exact}_{\mathrm{MIP}}
\]
가 된다. 또한 TS restricted 문제의 feasible solution 비용 $z^{R}_{\mathrm{MIP}}$가
같은 objective 기준으로 주어졌다면,
\[
  z^{R}_{\mathrm{MIP}} - z^{LB}_{\mathrm{MIP}}
\]
를 free-TS exact optimum에 대한 MIP objective gap의 상한으로 해석할 수
있다.

\begin{center}
\begin{tabular}{@{}p{0.24\linewidth}p{0.31\linewidth}p{0.31\linewidth}@{}}
\toprule
구분 & 기존 TS time restriction MCF & \texttt{mcf\_v6\_lb} \\
\midrule
정확한 원 문제와의 관계
& 정확한 TS timing feasible set에 추가 제약을 부과한다.
& 정확한 TS timing feasible set을 포함하도록 항해 가능성을 완화한다. \\
\addlinespace
feasible set 관계
& 실제 선박 최적해가 존재하는 경우 exact optimum의 상계 후보를 제공한다.
& no-virtual optimality 및 lower-envelope 비용 가정하에서 exact optimum의 하계를 제공한다. \\
\addlinespace
minimization 최적값 관계
& 같은 objective 기준에서 $z^{exact}_{\mathrm{MIP}} \le z^R_{\mathrm{MIP}}$
& 같은 objective 기준에서 $z^{LB}_{\mathrm{MIP}} \le z^{exact}_{\mathrm{MIP}}$ \\
\addlinespace
해석
& 실행 가능한 restricted schedule 비용을 제공한다.
& 실행 가능성을 보장하지 않을 수 있는 lower bound 비용을 제공한다. \\
\bottomrule
\end{tabular}
\end{center}

\section{공통 구조}

본 모형은 \texttt{mcf\_v5}와 같은 MIP 구조를 사용한다. 선박별 source-to-sink
flow, service arc coverage, position activation, arc capacity, TS unit activation,
TS in/out balance, dry dock category/linking 제약은 그대로 유지한다.

차이는 network construction 단계에서 TS 관련 노드의 시간과 이 시간을 사용하는
\texttt{SailArc} feasibility 계산에만 있다. 같은 항구를 잇는 경우에는 기존과
같이 항해가 필요 없는 일반 \texttt{Arc}로 생성한다.

현재 구현의 기본 설정에서는 가상 선박 commodity $v_0$를 최적화 모형에서 제외한다.
즉 기본 모형의 commodity 집합은 실제 선박 집합 $\mathcal{V}^{R}$뿐이며, $v_0$의
source node, target 종료 아크, wrap arc, flow 변수 및 virtual path 복원은 생성하지
않는다. 이 설정은 코드의 \texttt{INCLUDE\_VIRTUAL\_VESSEL = False}에 대응한다.
비교 실험을 위해 해당 옵션을 켜는 경우에만 \texttt{mcf\_v5}와 같은 $v_0$ commodity를
추가한다.

\section{가상 선박 옵션}

\texttt{INCLUDE\_VIRTUAL\_VESSEL = True}로 설정하면 $v_0$를 포함한 모형을 구성할 수
있다. 이 경우 $v_0$는 실제 재배치 선박이 아니라 service coverage를 위한 가상
commodity이므로, 서로 다른 lane-proforma-position에 속한 두 노드 그룹을 직접 잇는
base arc는 $v_0$의 아크 집합에서 제외한다. 이 제한은 서로 다른 항구 사이의
\texttt{SailArc}뿐 아니라, 두 노드 그룹의 항구가 같아 네트워크 construction 단계에서
일반 \texttt{Arc}로 생성된 same-port cross-lane arc에도 적용된다. $v_0$는 source에서
각 lane-proforma-position의 inbound nodes로 직접 phase-in할 수 있으므로, cross-lane
continuation arc가 없어도 가상 coverage 경로의 feasibility는 유지된다.

또한 $v_0$가 사용하는 base arc의 한쪽 끝이라도 TS 관련 노드
\[
  \texttt{TSIB},\quad \texttt{TSIA},\quad \texttt{TSOB},\quad \texttt{TSOA}
\]
에 해당하면 해당 arc의 $v_0$ opportunity cost는 0으로 둔다. \texttt{mcf\_v6\_lb}는
lower bound 생성을 위해 TS node time을 완화하므로, TS 관련 노드의 \texttt{in time}과
\texttt{out time}이 경제적 체류시간을 나타내지 않을 수 있다. 특히 완화된
\texttt{TSIA} 또는 \texttt{TSOB}에서는 내부 interval이 뒤집힐 수 있으므로, TS 관련
arc에 $v_0$ opportunity cost를 부과하면 인위적인 음수 objective coefficient가
생길 수 있다. 따라서 $v_0$가 켜진 경우에도 TS 관련 economic effect는 arc cost가
아니라 TS activation 변수의 비용으로만 반영한다.

\section{TS 시간 완화}

정규 port stay $g$의 pilot-in 시작 시각을 $t^g_s$, pilot-out 종료 시각을
$t^g_e$라고 한다. 같은 lane-position 안에서 $g$의 바로 이전 port stay를 $p(g)$,
바로 다음 port stay를 $n(g)$라고 한다. 두 port stay $i,j$의 항구 사이 거리를
$d_{ij}$라고 하고, 시간 단위가 hour일 때 $d_{ij}/20$은 20 knot로 항해할 때의
최소 항해 시간이다.

기존 MCF 계열 모형과 동일하게 환적 화물의 하역 시간과 선적 시간은 각각 6시간으로
둔다. 이전 또는 다음 port stay가 존재하지 않아 아래 식을 정의할 수 없는 TS unit은
exact feasible TS unit으로 보지 않는다.

\begin{center}
\renewcommand{\arraystretch}{1.2}
\begin{tabularx}{\textwidth}{@{}p{0.16\textwidth}XX@{}}
\toprule
노드 & 기존 TS time restriction & \texttt{mcf\_v6\_lb} relaxation \\
\midrule
\texttt{PLTI}$(g)$
& in: $t^g_s$, out: $t^g_s$
& in: $t^g_s$, out: $t^g_s$ \\
\addlinespace
\texttt{PLTO}$(g)$
& in: $t^g_e$, out: $t^g_e$
& in: $t^g_e$, out: $t^g_e$ \\
\addlinespace
\texttt{TSIB}$(g)$
& in: $t^g_s-6\mathrm{h}$, out: $t^g_s$
& in: $t^g_s-6\mathrm{h}$, out: $t^g_s$ \\
\addlinespace
\texttt{TSOB}$(g)$
& in: $t^g_s-36\mathrm{h}$, out: $t^g_s-30\mathrm{h}$
& in: $t^g_s-36\mathrm{h}$,
  out:
  $\begin{aligned}[t]
    \min\{&t^{p(g)}_e+d_{p(g),g}/20+6\mathrm{h},\\
           &t^g_s-30\mathrm{h}\}
  \end{aligned}$ \\
\addlinespace
\texttt{TSOA}$(g)$
& in: $t^g_e$, out: $t^g_e+6\mathrm{h}$
& in: $t^g_e$, out: $t^g_e+6\mathrm{h}$ \\
\addlinespace
\texttt{TSIA}$(g)$
& in: $t^g_e+30\mathrm{h}$, out: $t^g_e+36\mathrm{h}$
& in:
  $\begin{aligned}[t]
    \max\{&t^{n(g)}_s-d_{g,n(g)}/20-6\mathrm{h},\\
           &t^g_e+30\mathrm{h}\}
  \end{aligned}$,
  out: $t^g_e+36\mathrm{h}$ \\
\bottomrule
\end{tabularx}
\end{center}

\texttt{TSIA}의 in time은 고정 TS 시각 모형에서 쓰는 $t^g_e+30\mathrm{h}$보다
빠르지 않다. \texttt{TSOB}의 out time은 고정 TS 시각 모형에서 쓰는
$t^g_s-30\mathrm{h}$보다 늦지 않다. 따라서 변경된 TS 시간은 기존 고정 TS 시각
모형에서 가능한 \texttt{SailArc}의 항해 가능 시간을 줄이지 않는다.

\section{TS Activation Cost Lower Envelope}

TS unit $k$의 activation cost는 relaxed TS node의 대표 시각으로 계산하지 않는다.
TS-in node를 $\mathrm{TSI}_k$, TS-out node를 $\mathrm{TSO}_k$라 하고 각 node의
in/out time을 $\tau^{in}$, $\tau^{out}$으로 쓴다. 환적 하역 시작 시각의 가능한
month를 보수적으로 포함하기 위해 다음 window를 사용한다.
\[
  s_k^0 =
  \min\left\{
    \tau^{in}_{\mathrm{TSI}_k},
    \tau^{out}_{\mathrm{TSI}_k},
    \tau^{in}_{\mathrm{TSO}_k},
    \tau^{out}_{\mathrm{TSO}_k}-6\mathrm{h}
  \right\},
\]
\[
  e_k^0 =
  \max\left\{
    \tau^{in}_{\mathrm{TSI}_k},
    \tau^{out}_{\mathrm{TSI}_k},
    \tau^{in}_{\mathrm{TSO}_k},
    \tau^{out}_{\mathrm{TSO}_k}
  \right\}.
\]
Planning horizon을 $[H_s,H_e]$라 하면
\[
  s_k = \max\{s_k^0,H_s\},\qquad
  e_k = \min\{e_k^0,H_e\}
\]
로 둔다. $s_k>e_k$이면 해당 TS unit의 cost window가 planning horizon과 교차하지
않으므로 네트워크 생성 데이터가 유효하지 않은 것으로 보고 에러를 발생시킨다.

정확 문제에서 해당 TS unit의 환적 하역 시작 시각을 $\alpha$라고 할 때,
\[
  c^{LB}_k
  =
  \min_{\alpha\in [s_k,e_k]}
  \mathrm{TSPrice}\bigl(\mathrm{lane}(k),\mathrm{port}(k),
  \mathrm{yearmonth}(\alpha)\bigr)
\]
로 둔다. 구현은 $[s_k,e_k]$의 시작과 끝 사이에 등장하는 모든 year-month를 포함한다.
어떤 year-month의 TS price가 입력 데이터에 없으면 임의 fallback을 사용하지 않고
에러를 발생시킨다.

\section{TS-Related Sailing Cost Lower Envelope}

완화된 \texttt{TSIA}, \texttt{TSOB} node time은 항해 feasibility를 완화하기 위한
인공 시간이다. 따라서 이 시각의 year-month를 그대로 bunker price month로 사용하면
exact 가능한 항해 비용보다 큰 비용이 부과될 수 있다. 이를 피하기 위해 한쪽 끝이
TS 관련 노드인 \texttt{SailArc} 또는 \texttt{HorizonSailArc}의 항해 비용은 입력
데이터에 존재하는 bunker price year-month 전체에 대해 계산한 비용의 최솟값으로
둔다. 이는 month attribution에 대한 lower envelope이며, TS timing relaxation의
경제 비용이 exact 비용보다 커지는 것을 방지하기 위한 보수적 처리이다.

각 year-month에서의 bunker price 계산은 \texttt{mcf\_v5}의 MIP objective와 같은
가격 귀속 규칙을 사용한다. 즉 해당 year-month, lane, fuel type의 lane-specific
price가 있으면 그 값을 쓰고, 없으면 같은 year-month와 fuel type에 대해 입력 데이터에
존재하는 가격들의 산술평균을 쓴다. 따라서 TS 관련 항해의 lower-envelope 비용은
``\texttt{mcf\_v5}와 동일한 bunker price convention으로 각 가능한 month의 항해
비용을 계산한 뒤 그 최솟값을 취한 값''이다. 이 lower bound statement는 exact 문제도
같은 bunker price convention을 사용한다는 조건하에서 해석한다.

\section{SailArc Feasibility}

서로 다른 항구를 잇는 후보 아크 $a=(u,v)$에 대해 거리는 기존과 같이
$d(u,v)$를 사용한다. 항해 가능 시간은 변경된 노드 시간을 사용하여
\[
  s(u,v)=\tau^{in}_v-\tau^{out}_u
\]
로 계산한다. 후보 아크는
\[
  s(u,v)\ge 0,\qquad \frac{d(u,v)}{s(u,v)} \le 20
\]
을 만족하는 경우에만 \texttt{SailArc}로 생성한다. 같은 항구를 잇는 후보 아크는
항해가 필요 없으므로 기존과 같이 일반 \texttt{Arc}로 생성한다.

\section{실행 결과와 비교 기준}

\texttt{mcf\_v6\_lb}는 실행 결과로 feasible vessel schedule을 반환하지 않는다.
구현상 알고리즘의 반환 객체는 \texttt{LowerBoundResult}이며, orchestration layer는
이 결과를 solution validation, OCAM evaluation, postprocessing 대상으로 취급하지
않는다. 저장되는 값은 lower bound 값과 최적화 모형의 보조 정보이다. 현재 저장
형식에는 lower bound 값, Gurobi objective bound, MIP gap, solver status, solution
count, model name이 포함된다.

비교 시 absolute gap은
\[
  z^{v4}-z^{LB}
\]
로 계산하고, relative gap은
\[
  \frac{z^{v4}-z^{LB}}{|z^{v4}|}
\]
로 계산한다. \texttt{mcf\_v6\_lb} 모형이 \texttt{OPTIMAL} 상태로 종료되고 저장된
MIP gap이 $0$이면 해당 lower bound는 relaxation 모형의 최적값이다.
이 값은 feasible schedule cost가 아니므로 \texttt{mcf\_v5} solution의 대체 결과가
아니라, restricted feasible solution이 같은 MIP objective 기준의 free-TS
exact optimum에서 얼마나 멀 수 있는지를 평가하기 위한 기준값으로 사용한다.

\end{document}
